%! Justifying a Claim Based on a Confidence Interval for a Population Proportion — Unit 6, Lesson 6.3 — Exit Ticket
\documentclass[10pt]{article}
\usepackage{apstats-article}
\usepackage{apstats-boxes}
\newcommand{\TallMath}[1]{$\displaystyle #1\rule[-1.4em]{0pt}{3.2em}$}

\begin{document}

\pageheader{Unit 6, Lesson 6.3}{Exit Ticket}
\namedateperiod

\vspace{0.15in}

A random sample of 220 high school students found that 143 prefer to study with music
playing. A 90\% confidence interval for the true proportion is $(0.592, 0.709)$.

\begin{enumerate}
  \item Write a complete interpretation of this confidence interval.
        \writelines{2}

  \item A researcher claims that exactly 65\% of high school students prefer to study with
        music. Is this claim plausible based on the interval? Explain.
        \writelines{2}

  \item A school principal claims that \emph{more than} 70\% of students prefer studying
        with music. Does the CI support or refute this claim? Explain.
        \writelines{2}

  \item Explain in one sentence what ``90\% confident'' means about the procedure, not about
        this particular interval.
        \writelines{2}
\end{enumerate}

\end{document}
